\section{引言}
\label{sec:1}
尽管格点正规化为场论提供了非常强有力的非微扰表述，但遗憾的是，它与基本的
整体对称性常常不相容。最著名的例子是手征对称性~\cite{Nielsen:1980rz,%
Nielsen:1981xu}；超对称性是另一个臭名昭著的例子~\cite{Dondi:1976tx}，不用
说，平移不变性也是如此。当一种正规化在某个对称性下不是不变的时候，要构造
相应的 Noether 流并不直接了当——该流应当守恒，并通过 Ward--Takahashi
（WT）关系生成对称性变换。这使得在坚实的基础上测量与 Noether 流相关的物
理量十分困难。为解决这个问题，人们可以设想至少三种可能的途径。

第一种途径是理想化的：找到一种实现（格点修正形式的）所求对称性的格点表
述。如果这样的表述唾手可得，那么用标准的 Noether 方法就可以容易地得到相
应的 Noether 流。这方面最成功的例子是格点手征对称性~\cite{Kaplan:1992bt,%
Neuberger:1997fp,Hasenfratz:1997ft,Hasenfratz:1998ri,Neuberger:1998wv,%
Hasenfratz:1998jp,Luscher:1998pqa,Niedermayer:1998bi}，它可以由满足
Ginsparg--Wilson 关系~\cite{Ginsparg:1981bj} 的格点狄拉克算符来定义。虽然
这无疑是一条理想的途径，但这样的理想表述似乎并不总能得到，特别是对时空对
称性（关于超对称性的 no-go 定理，例如见文献~\cite{Kato:2008sp}）。

第二种途径是通过调节系数来构造 Noether 流，这些系数出现在按格点对称性可
与 Noether 流混合的算符的线性组合之中。\footnote{这里，我们假设已经完成
了把裸参数细调到目标（对称）理论的工作。}例如，对于能量--动量张量——与平
移不变性以及转动和共形对称性相联系的 Noether
流~\cite{Callan:1970ze,Coleman:1970je}——可以通过调节维数~$4$ 算符线性组
合中的系数，构造出守恒的格点能量--动量
张量~\cite{Caracciolo:1988hc,Caracciolo:1989pt}\footnote{基于
$\mathcal{N}=1$ 超对称性的一个稍有不同的途径见文
献~\cite{Suzuki:2013gi,Suzuki:2012wx}。}；而能量--动量张量的整体归一化必
须用其他方式确定。\footnote{为此也许可以采用“流代数”，正如轴矢流的情形
~\cite{Bochicchio:1985xa}。}虽然当能量--动量张量处于“孤立”状态时——即当
能量--动量张量与其他复合算符分离开时（如壳上矩阵元的情形）——这一方法在
原则上已经足够，但先验地并不清楚：当能量--动量张量与其他复合算符在位置空
间重合时，是否能够控制可能参与贡献的高维算符带来的不确定性。这意味着，我
们并不显然知道：用上述方法构造的能量--动量张量通过 WT 关系作用在算符上
时，是否生成正确归一化的平移（以及转动和共形变换）。（如果能量--动量张量
生成正确归一化的平移，则可以保证~\cite{Fujikawa:1980rc}（另见文
献~\cite{Fujikawa:2004cx} 第 7.3 节），迹反常或共形反常~\cite{Crewther:1972kn,%
Chanowitz:1972vd} 正比于重正化群函数~\cite{Adler:1976zt,Nielsen:1977sy,%
Collins:1976yq}。）

第三种可能的途径是利用某种紫外（UV）有限的量。由于这样的量必定与所采用的
正规化无关（在取掉调节子的极限下），就出现了这样一种可能性：可以把格点正
规化与某种保持所求对称性的其他正规化联系起来。这一方法论可见于例如文
献~\cite{Luscher:2004fu}（另见文献~\cite{Luscher:2010ik}），其中导出了拓
扑磁化率的紫外有限表示。虽然该表示本身的推导依赖于保持手征对称性的格点正
规化~\cite{Kaplan:1992bt,Neuberger:1997fp,Hasenfratz:1997ft,%
Hasenfratz:1998ri,Neuberger:1998wv,Hasenfratz:1998jp,Luscher:1998pqa,%
Niedermayer:1998bi}，但由于它必定与正规化无关，人们可以用任何正规化（例
如 Wilson 费米子~\cite{Wilson:1975hf}）来计算这个表示。

本文以纯杨--米尔斯理论为例，针对能量--动量张量考虑上述第三种途径。为此，
我们利用所谓的杨--米尔斯梯度流（在格点规范理论的语境中即 Wilson 流），其
在格点规范理论中的有用性最近已被揭示~\cite{Luscher:2010iy,Luscher:2010we,%
Luscher:2011bx,Borsanyi:2012zs,Borsanyi:2012zr,Fodor:2012td,Fodor:2012qh,%
Fritzsch:2013je,Luscher:2013cpa}。杨--米尔斯梯度流的一个显著特点是其稳健
的 UV 有限性~\cite{Luscher:2011bx}。更确切地说，由梯度流产生的规范场的任
意乘积，对正的流时间~$t$，在标准重正化之下是 UV 有限的。而且，即使规范场
的某些位置相互重合，这样的乘积也\emph{不\/}表现出任何奇异性。这种 UV 有
限性的基本机制是：流方程是一类扩散方程，而动量空间中的演化算符~$\sim
e^{-tk^2}$ 在~$t>0$ 时起到 UV 调节子的作用。梯度流的这一性质意味着，正流
时间处规范场局域乘积的定义与正规化无关。在我们当前的语境中，有希望把格点
正规化得到的量与维数正规化得到的量联系起来——对于后者，平移不变性\emph{是\/}
明显成立的。

另一方面，如文献~\cite{Luscher:2011bx} 所指出的，正流时间处规范场的局域乘
积可以用原来规范理论的重正化局域算符展开，展开系数有限。这些系数满足一定
的重正化群方程；结合量纲分析，它们给出系数作为流时间函数的信息。由于渐近
自由，进而可以用微扰论求出这些系数在小流时间下的渐近行为。

利用上述梯度流的性质，可以得到一个公式，它把某些规范不变局域乘积的小流时
间行为与由维数正规化定义的能量--动量张量联系起来。由于前者可以用格点正规
化的 Wilson 流计算~\cite{Luscher:2010iy,Luscher:2010we,Luscher:2011bx,%
Borsanyi:2012zs,Borsanyi:2012zr,Fodor:2012td,Fodor:2012qh,Fritzsch:2013je,%
Luscher:2013cpa}，而后者守恒并在复合算符上生成正确归一化的平移，我们的公
式提供了一种可能的方法：用蒙特卡罗模拟来计算正确归一化的守恒能量--动量张
量的关联函数。

除非另有说明，本文沿用文献~\cite{Luscher:2011bx} 的记号约定。
